% ---------------------------------------------------------------------------
% Study report · Predicting the inclination from the legacy sensor set
%
% Build:  pdflatex inclination_prediction_legacy_report.tex  (twice, for the ToC)
% Figures are read from ../outputs/ and are produced by the paired notebook.
% ---------------------------------------------------------------------------
\documentclass[11pt,a4paper]{article}

\newcommand{\runninghead}{Legacy-era prediction · station 02 · 2018--2025}

% ---------------------------------------------------------------------------
% reportstyle.tex — shared preamble for the study reports under studies/.
%
% Kept deliberately inside the package set shipped with TeX Live "basic":
% no tcolorbox, no siunitx, no enumitem, no titlesec. Everything below
% compiles with a stock pdflatex installation.
%
% Usage, from a report directory one level below a study folder:
%     % ---------------------------------------------------------------------------
% reportstyle.tex — shared preamble for the study reports under studies/.
%
% Kept deliberately inside the package set shipped with TeX Live "basic":
% no tcolorbox, no siunitx, no enumitem, no titlesec. Everything below
% compiles with a stock pdflatex installation.
%
% Usage, from a report directory one level below a study folder:
%     % ---------------------------------------------------------------------------
% reportstyle.tex — shared preamble for the study reports under studies/.
%
% Kept deliberately inside the package set shipped with TeX Live "basic":
% no tcolorbox, no siunitx, no enumitem, no titlesec. Everything below
% compiles with a stock pdflatex installation.
%
% Usage, from a report directory one level below a study folder:
%     \input{../../_shared/reportstyle.tex}
%     \graphicspath{{../outputs/}}
% ---------------------------------------------------------------------------

\usepackage[T1]{fontenc}
\usepackage[utf8]{inputenc}
\usepackage{textcomp}
\usepackage{lmodern}
\usepackage{microtype}
\usepackage[margin=2.4cm,headheight=15pt]{geometry}
\usepackage{graphicx}
\usepackage{booktabs}
\usepackage{tabularx}
\usepackage{array}
\usepackage{amsmath}
\usepackage{xcolor}
\usepackage[font=small,labelfont=bf,justification=justified]{caption}
\usepackage{float}
\usepackage{fancyhdr}
\usepackage[hidelinks]{hyperref}

% --- Palette ---------------------------------------------------------------
\definecolor{rulegrey}{HTML}{B9C0C7}
\definecolor{fillgrey}{HTML}{F2F4F6}
\definecolor{failred}{HTML}{A5202B}
\definecolor{passgreen}{HTML}{1B6B3A}
\definecolor{accent}{HTML}{1F4E79}

% --- Semantic shorthands ---------------------------------------------------
\newcommand{\degC}{\textdegree C}
\newcommand{\mdegC}{mdeg/\textdegree C}
\newcommand{\fail}{\textcolor{failred}{\textbf{fail}}}
\newcommand{\pass}{\textcolor{passgreen}{\textbf{pass}}}
\newcommand{\worse}{\textcolor{failred}{worse}}
\newcommand{\better}{\textcolor{passgreen}{better}}

% --- Highlighted panel -----------------------------------------------------
% #1 = heading, #2 = body
\newcommand{\panel}[2]{%
  \begin{center}%
  \setlength{\fboxsep}{9pt}%
  \fcolorbox{rulegrey}{fillgrey}{%
    \begin{minipage}{0.93\textwidth}%
      {\bfseries\color{accent}#1}\par\medskip
      #2
    \end{minipage}}%
  \end{center}}

% --- Numeric column: right aligned, fixed width ----------------------------
\newcolumntype{R}[1]{>{\raggedleft\arraybackslash}p{#1}}
\newcolumntype{L}[1]{>{\raggedright\arraybackslash}p{#1}}
\newcolumntype{Y}{>{\raggedright\arraybackslash}X}

% --- Table look ------------------------------------------------------------
\renewcommand{\arraystretch}{1.15}
\setlength{\tabcolsep}{6pt}

% --- Headers ---------------------------------------------------------------
\pagestyle{fancy}
\fancyhf{}
\fancyhead[L]{\small\itshape\runninghead}
\fancyhead[R]{\small\thepage}
\renewcommand{\headrulewidth}{0.4pt}

% --- Section spacing without titlesec --------------------------------------
\makeatletter
\renewcommand\section{\@startsection{section}{1}{\z@}%
  {-3.2ex \@plus -1ex \@minus -.2ex}{2.0ex \@plus.2ex}%
  {\normalfont\Large\bfseries\color{accent}}}
\renewcommand\subsection{\@startsection{subsection}{2}{\z@}%
  {-2.6ex \@plus -1ex \@minus -.2ex}{1.4ex \@plus.2ex}%
  {\normalfont\large\bfseries}}
\makeatother

% --- Figure defaults -------------------------------------------------------
\setkeys{Gin}{width=\linewidth}
\captionsetup[figure]{skip=6pt}

%     \graphicspath{{../outputs/}}
% ---------------------------------------------------------------------------

\usepackage[T1]{fontenc}
\usepackage[utf8]{inputenc}
\usepackage{textcomp}
\usepackage{lmodern}
\usepackage{microtype}
\usepackage[margin=2.4cm,headheight=15pt]{geometry}
\usepackage{graphicx}
\usepackage{booktabs}
\usepackage{tabularx}
\usepackage{array}
\usepackage{amsmath}
\usepackage{xcolor}
\usepackage[font=small,labelfont=bf,justification=justified]{caption}
\usepackage{float}
\usepackage{fancyhdr}
\usepackage[hidelinks]{hyperref}

% --- Palette ---------------------------------------------------------------
\definecolor{rulegrey}{HTML}{B9C0C7}
\definecolor{fillgrey}{HTML}{F2F4F6}
\definecolor{failred}{HTML}{A5202B}
\definecolor{passgreen}{HTML}{1B6B3A}
\definecolor{accent}{HTML}{1F4E79}

% --- Semantic shorthands ---------------------------------------------------
\newcommand{\degC}{\textdegree C}
\newcommand{\mdegC}{mdeg/\textdegree C}
\newcommand{\fail}{\textcolor{failred}{\textbf{fail}}}
\newcommand{\pass}{\textcolor{passgreen}{\textbf{pass}}}
\newcommand{\worse}{\textcolor{failred}{worse}}
\newcommand{\better}{\textcolor{passgreen}{better}}

% --- Highlighted panel -----------------------------------------------------
% #1 = heading, #2 = body
\newcommand{\panel}[2]{%
  \begin{center}%
  \setlength{\fboxsep}{9pt}%
  \fcolorbox{rulegrey}{fillgrey}{%
    \begin{minipage}{0.93\textwidth}%
      {\bfseries\color{accent}#1}\par\medskip
      #2
    \end{minipage}}%
  \end{center}}

% --- Numeric column: right aligned, fixed width ----------------------------
\newcolumntype{R}[1]{>{\raggedleft\arraybackslash}p{#1}}
\newcolumntype{L}[1]{>{\raggedright\arraybackslash}p{#1}}
\newcolumntype{Y}{>{\raggedright\arraybackslash}X}

% --- Table look ------------------------------------------------------------
\renewcommand{\arraystretch}{1.15}
\setlength{\tabcolsep}{6pt}

% --- Headers ---------------------------------------------------------------
\pagestyle{fancy}
\fancyhf{}
\fancyhead[L]{\small\itshape\runninghead}
\fancyhead[R]{\small\thepage}
\renewcommand{\headrulewidth}{0.4pt}

% --- Section spacing without titlesec --------------------------------------
\makeatletter
\renewcommand\section{\@startsection{section}{1}{\z@}%
  {-3.2ex \@plus -1ex \@minus -.2ex}{2.0ex \@plus.2ex}%
  {\normalfont\Large\bfseries\color{accent}}}
\renewcommand\subsection{\@startsection{subsection}{2}{\z@}%
  {-2.6ex \@plus -1ex \@minus -.2ex}{1.4ex \@plus.2ex}%
  {\normalfont\large\bfseries}}
\makeatother

% --- Figure defaults -------------------------------------------------------
\setkeys{Gin}{width=\linewidth}
\captionsetup[figure]{skip=6pt}

%     \graphicspath{{../outputs/}}
% ---------------------------------------------------------------------------

\usepackage[T1]{fontenc}
\usepackage[utf8]{inputenc}
\usepackage{textcomp}
\usepackage{lmodern}
\usepackage{microtype}
\usepackage[margin=2.4cm,headheight=15pt]{geometry}
\usepackage{graphicx}
\usepackage{booktabs}
\usepackage{tabularx}
\usepackage{array}
\usepackage{amsmath}
\usepackage{xcolor}
\usepackage[font=small,labelfont=bf,justification=justified]{caption}
\usepackage{float}
\usepackage{fancyhdr}
\usepackage[hidelinks]{hyperref}

% --- Palette ---------------------------------------------------------------
\definecolor{rulegrey}{HTML}{B9C0C7}
\definecolor{fillgrey}{HTML}{F2F4F6}
\definecolor{failred}{HTML}{A5202B}
\definecolor{passgreen}{HTML}{1B6B3A}
\definecolor{accent}{HTML}{1F4E79}

% --- Semantic shorthands ---------------------------------------------------
\newcommand{\degC}{\textdegree C}
\newcommand{\mdegC}{mdeg/\textdegree C}
\newcommand{\fail}{\textcolor{failred}{\textbf{fail}}}
\newcommand{\pass}{\textcolor{passgreen}{\textbf{pass}}}
\newcommand{\worse}{\textcolor{failred}{worse}}
\newcommand{\better}{\textcolor{passgreen}{better}}

% --- Highlighted panel -----------------------------------------------------
% #1 = heading, #2 = body
\newcommand{\panel}[2]{%
  \begin{center}%
  \setlength{\fboxsep}{9pt}%
  \fcolorbox{rulegrey}{fillgrey}{%
    \begin{minipage}{0.93\textwidth}%
      {\bfseries\color{accent}#1}\par\medskip
      #2
    \end{minipage}}%
  \end{center}}

% --- Numeric column: right aligned, fixed width ----------------------------
\newcolumntype{R}[1]{>{\raggedleft\arraybackslash}p{#1}}
\newcolumntype{L}[1]{>{\raggedright\arraybackslash}p{#1}}
\newcolumntype{Y}{>{\raggedright\arraybackslash}X}

% --- Table look ------------------------------------------------------------
\renewcommand{\arraystretch}{1.15}
\setlength{\tabcolsep}{6pt}

% --- Headers ---------------------------------------------------------------
\pagestyle{fancy}
\fancyhf{}
\fancyhead[L]{\small\itshape\runninghead}
\fancyhead[R]{\small\thepage}
\renewcommand{\headrulewidth}{0.4pt}

% --- Section spacing without titlesec --------------------------------------
\makeatletter
\renewcommand\section{\@startsection{section}{1}{\z@}%
  {-3.2ex \@plus -1ex \@minus -.2ex}{2.0ex \@plus.2ex}%
  {\normalfont\Large\bfseries\color{accent}}}
\renewcommand\subsection{\@startsection{subsection}{2}{\z@}%
  {-2.6ex \@plus -1ex \@minus -.2ex}{1.4ex \@plus.2ex}%
  {\normalfont\large\bfseries}}
\makeatother

% --- Figure defaults -------------------------------------------------------
\setkeys{Gin}{width=\linewidth}
\captionsetup[figure]{skip=6pt}

\graphicspath{{../outputs/}}

\title{\vspace{-1.2cm}\textbf{Predicting the inclination of the Gubbio wall
from the legacy sensor set}\\[2mm]
\large Station 02, 26 July 2018 to 20 February 2025, hourly grid}
\author{Study report · \texttt{studies/inclination\_prediction\_legacy/}}
\date{12 August 2026}

\begin{document}
\maketitle
\thispagestyle{fancy}

\begin{abstract}
\noindent
The 14-column network ran at station 02 for six and a half years before the
February 2025 changeover. It carries air temperature, relative humidity, battery
voltage and inclination --- no solar radiation, no wall temperature.

This report asks two families of question. The first is what the extra length
buys, and is specific to this era: which periodicities are present, over which
stretch they should be decomposed, and how missing values should be filled. The
second is the set the extended-sensor study asked --- best single predictor, best
combination and how to diagnose it, how far ahead and with what uncertainty ---
repeated here with that study's machinery unchanged, so that any difference in
result is a difference in the data and not in the method.

\textbf{The periodicities are real and identified} without imputing anything
first: an annual term at 370\,d, a semi-annual term at 183\,d, and a diurnal term
at 23.93\,h with its second and third harmonics. A window-function control shows
the apparent 971-day feature to be an artefact of the outage pattern.
\textbf{The semi-annual term is the annual second harmonic}, not an independent
180-day process: frequency alone can never separate the two on a record of this
length, but they are phase locked to within 28.5\textdegree.
\textbf{The components do not hold still} --- the annual amplitude ranges from
30.6 to 53.3\,mdeg --- and that is what makes the seasonal model useless for
prediction. \textbf{Almost all of the forecast skill is autoregressive}: one hour
ahead, 47.7 of the 48.8\,mdeg climatological error is removed by the signal's own
recent history, the seasonal terms remove 0.1, and the drivers remove nothing.

The ported questions answer sharply and negatively. Air temperature is the best
single predictor on the whole era at $R^2 = 0.073$; on the unbroken block
\emph{no driver beats a constant at all}. The two-driver combination improves on
air temperature by 1.3\,\% against a fold spread twenty-three times that size,
and out of period both models roughly double their error. Two warnings for the
framework follow: extrapolating a fitted linear trend costs \textbf{85--93\,mdeg
at every horizon}, and only about a tenth of the missing hours can be defensibly
filled.

\textbf{A seventh question was raised by those answers and is addressed in
Section~\ref{sec:usable}: how much of the record can be used at all.} Everything
above was computed under an assumption stated in Section~\ref{sec:record} and
never tested --- that autoregressive work must run on one unbroken block. It need
not. An autoregressive model consumes windows, not a continuous series, and
fitting across every complete stretch of the era yields \textbf{6.57 annual
cycles and 2.26 times the training windows for 1.5\,\% fabricated data}. The
consequence is large: the deployable configuration's prediction intervals fall
from 164--185\,mdeg, which excluded nothing, to \textbf{12--23\,mdeg}, and its
one-hour error by 44\,\%, with no change to the model. The same comparison tests
a rule the companion framework asserts --- that one to two annual cycles suffice
for a yearly component --- and finds it half right: the longer span corrects the
annual amplitude but lets NeuralProphet's trend absorb \emph{more} of the
low-frequency variance, 53.8\,\% against 39.9.
\end{abstract}

\tableofcontents

% ===========================================================================
\section{What this study asks, and what it assumes}

\subsection{Six questions, three of them ported}

\textbf{Specific to the long record.}
\begin{enumerate}
\item Which periodicities are present in the compensated inclination?
\item Over which stretch should they be decomposed, and do they persist?
\item How should missing values be filled, and up to what gap length is that
defensible?
\end{enumerate}

\textbf{Ported from \texttt{studies/inclination\_prediction/}, for
comparability.}
\begin{enumerate}
\setcounter{enumi}{3}
\item Can one sensor channel predict the inclination, and which is best?
\item Can a combination beat the best single channel, and how is the best
combination diagnosed?
\item How far ahead can NeuralProphet forecast, with what uncertainty --- and
what is the error at each horizon actually made of?
\end{enumerate}

Question 6 merges the sibling study's horizon question with this study's error
budget. That is the natural place for it: a horizon curve says how large the
error is, and the budget says which component produced it.

\subsection{Premises, imposed rather than tested}

\begin{center}
\begin{tabularx}{\textwidth}{L{4.6cm}Y}
\toprule
\textbf{Premise} & \textbf{Consequence for this report} \\
\midrule
The logged \texttt{I} channel is an inclination in millidegrees &
The 2500 offset is added back. In this era the logged values run 2105--2300, so
read as millidegrees they lie \textbf{outside} the $\pm$2000\,mdeg the datasheet
certifies. Accepted by instruction; the compensation normalises the series to
start at zero, so nothing downstream depends on it. \\
The documented compensation is correct &
\texttt{inc\_comp} is the target of every model. Section~\ref{sec:sensitivity}
shows this is by far the most consequential of the three. \\
The analysis grid is one hour &
Matching every external proxy the pipeline aligns against. \\
\bottomrule
\end{tabularx}
\end{center}

\subsection{The sign contradiction, reproduced on a second instrument}
\label{sec:sign}

The site predicts a negative association between the inclination and any heating
driver: the valley face is the more sun-exposed, expands further, and tips the
wall towards the mountain. As in the extended-sensor study, that prediction is
tested on the \emph{raw} channel, before the compensation can manufacture the
expected sign.

\begin{center}
\small
\begin{tabular}{llR{1.4cm}R{1.7cm}R{1.4cm}R{1.7cm}}
\toprule
\textbf{Series} & \textbf{Driver} & \textbf{$r$ lev.} & \textbf{slope lev.} &
\textbf{$r$ band} & \textbf{slope band} \\
\midrule
Raw \texttt{inc}   & \texttt{tair} & $+0.241$ & $+1.173$ & $+0.957$ & $+2.379$ \\
                   & \texttt{rh}   & $-0.244$ & $-0.479$ & $-0.584$ & $-0.343$ \\
\texttt{inc\_comp} & \texttt{tair} & $-0.629$ & $-3.827$ & $-0.964$ & $-2.621$ \\
                   & \texttt{rh}   & $+0.301$ & $+0.738$ & $+0.595$ & $+0.381$ \\
\bottomrule
\end{tabular}
\end{center}

\noindent{\footnotesize Slopes in mdeg per \degC\ for temperature and per
percentage point for humidity; ``band'' is the diurnal band.\par}

\panel{The same positive slope on a different instrument}{%
The legacy inclinometer moves \emph{positively} with air temperature, at
$+2.379$\,\mdegC\ on the diurnal band. The extended-package instrument, installed
seven years later as a physical reinstallation sharing no baseline and no
hardware with this one, measured $+2.291$.

\medskip
Two independent instruments, agreeing to within 4\,\% on a slope that runs
opposite to the site's structural expectation. That is much harder to explain as
a per-instrument defect than as a systematic property of the instrument type or
of the sign convention itself, and it correspondingly weakens the thermal-drift
reading of \texttt{docs/raw-data-format.md} Section~7.5 in favour of an inverted
convention.

\medskip
This is new evidence and it does not settle the question, which still needs the
calibration sheet or the mounting orientation. It does mean the contradiction is
a property of the measurement chain rather than of one sensor.}

% ===========================================================================
\section{The record, and what it can support}
\label{sec:record}

\subsection{Loading and cleaning}

Of the 1926 legacy files, a third carry duplicate timestamps whose copies
disagree, an eighth mix decimal separators internally, and some carry records
dated to the previous day. Duplicates are merged field by field rather than
dropped, because where copies differ it is usually because one carries real
values where the other carries zeros: 19\,210 duplicated records were merged,
resolving 18\,473 field-level conflicts. Sentinels sent 64\,248 zeros to
\texttt{NaN}. The spike screen removed nothing --- unlike the current era, this
record has no acquisition-restart transient.

The result is 57\,625 hourly slots of which 39\,499 carry the inclinometer, that
is \textbf{68.5\,\% coverage}.

\begin{figure}[H]
\includegraphics{S4_F01_st02_overview}
\caption{The legacy era in full. The compensated inclination spans roughly
285\,mdeg against a diurnal amplitude of tens, and the outages are visible as the
blank stretches --- five of them longer than thirty days.}
\end{figure}

\begin{figure}[H]
\includegraphics{S4_F02_st02_target}
\caption{Raw and compensated inclination with the correction that separates them.
As in the current era, the correction is not a small adjustment.}
\end{figure}

\subsection{Where the missing hours are}

\begin{center}
\small
\begin{tabular}{lR{1.6cm}R{2.0cm}R{2.4cm}}
\toprule
\textbf{Gap length} & \textbf{gaps} & \textbf{hours lost} &
\textbf{share of missing} \\
\midrule
$\leq$ 3 h   & 417 &   606 &  3.3\,\% \\
3--6 h       &  58 &   275 &  1.5\,\% \\
6--24 h      &  70 &   875 &  4.8\,\% \\
1--3 d       &   2 &    65 &  0.4\,\% \\
3--7 d       &   0 &     0 &  0.0\,\% \\
7--30 d      &   2 &   989 &  5.5\,\% \\
$>$ 30 d     &   5 & 15\,316 & \textbf{84.5\,\%} \\
\bottomrule
\end{tabular}
\end{center}

The shape of this table decides the imputation question before any method is
tried. There are 554 gaps, but 84.5\,\% of the lost hours sit in five of them.

\subsection{No block contains two annual cycles}
\label{sec:blocks}

\begin{center}
\small
\begin{tabular}{lR{2.2cm}R{1.6cm}R{1.8cm}R{1.6cm}}
\toprule
\textbf{Longest blocks, 72 h tolerance} & \textbf{start} & \textbf{days} &
\textbf{coverage} & \textbf{cycles} \\
\midrule
first  & 2020-07-31 & 617.2 & 96.8\,\% & \textbf{1.69} \\
second & 2019-05-11 & 390.1 & 97.7\,\% & 1.07 \\
third  & 2023-06-21 & 317.5 & 91.8\,\% & 0.87 \\
fourth & 2024-10-09 & 133.4 & 99.3\,\% & 0.37 \\
\bottomrule
\end{tabular}
\end{center}

\panel{This divides the study in two}{%
Even absorbing interruptions of up to three days, the longest unbroken stretch
holds \textbf{1.69} annual cycles. Nothing in this record qualifies on a
two-cycle rule.

\medskip
The annual term therefore cannot be identified from any single contiguous window.
It has to be fitted \emph{across} the gaps --- which a least-squares harmonic fit
does naturally and an autoregressive model cannot. So the spectrum and the
harmonic decomposition run on the \textbf{whole era}; the autoregressive work
runs on the \textbf{longest block}, 2020-07-31 to 2022-04-09; and the third
block, 2023-06-21 to 2024-05-04, is held back entirely for the out-of-period test
of Section~\ref{sec:holdout}.}

% ===========================================================================
\section{Question 1 --- the spectrum}

\subsection{Why not a periodogram}

A fifth of the record is missing and the missingness is structured into a handful
of long outages. An ordinary periodogram needs a complete grid, and filling those
holes before estimating the spectrum lets the filling method decide the answer: a
271-day hole interpolated linearly injects power at exactly the low frequencies
under investigation.

The Lomb--Scargle periodogram fits a sinusoid at each trial frequency by least
squares over whatever samples exist, so it consumes the gaps rather than
requiring them filled. \textbf{The spectrum is computed before any imputation.}

Two controls accompany it. The \textbf{window function} is the same periodogram
run on the observation mask alone: any peak appearing there is a property of when
the instrument was running. The \textbf{permutation null} shuffles values against
observation times, destroying periodicity while preserving the sampling pattern
exactly; two hundred replicates put the 95th and 99th percentiles at 0.0124 and
0.0140.

\subsection{What is there}

\begin{center}
\small
\begin{tabular}{lR{1.6cm}R{1.6cm}R{2.2cm}L{3.4cm}}
\toprule
\textbf{Period} & \textbf{power} & \textbf{window} & \textbf{above 99\,\%} &
\textbf{reading} \\
\midrule
370.1 d (1.01 yr) & \textbf{0.3995} & 0.0279 & yes & annual, genuine \\
971.0 d (2.66 yr) & 0.1178 & \textbf{0.1336} & yes & \fail\ outage artefact \\
544.6 d (1.49 yr) & 0.0538 & 0.0207 & yes & no physical reading \\
183.0 d (0.50 yr) & 0.0351 & 0.0211 & yes & semi-annual, genuine \\
113.9 d (0.31 yr) & 0.0278 & 0.0081 & yes & unanticipated, genuine \\
75.2 d            & 0.0052 & 0.0027 & no  & below the null \\
\bottomrule
\end{tabular}
\end{center}

\begin{figure}[H]
\includegraphics{S4_F03_st02_spectrum_low}
\caption{Low-frequency Lomb--Scargle spectrum on daily means, with the
permutation levels, and the window function beneath. The rise in the window
function beyond about 700 days is what disqualifies the 971-day feature.}
\end{figure}

\panel{The window-function control earns its place}{%
The second-largest peak in the data spectrum, at 971 days, has a
\emph{window-function power larger than its own} --- 0.1336 against 0.1178. It is
produced by the arrangement of the outages, not by the wall. Without this control
it would have been reported as a 2.7-year structural cycle, comfortably clearing
the permutation null.}

The three periodicities anticipated before the analysis --- one day, roughly six
months, one year --- are all present and all clear the null. A fourth, at
\textbf{113.9 days}, was not anticipated. It is well separated from the window
function, so it is not an outage artefact, and it has no obvious harmonic
relationship to the annual term. This report records it and does not explain it.

\subsection{The diurnal band}

\begin{center}
\small
\begin{tabular}{lR{1.8cm}L{6.2cm}}
\toprule
\textbf{Period} & \textbf{power} & \textbf{reading} \\
\midrule
23.93 h & \textbf{0.0407} & diurnal \\
7.28 d  & 0.0068 & \fail\ artefact of the 168\,h detrending filter \\
11.99 h & 0.0017 & second harmonic \\
8.00 h  & 0.0015 & third harmonic \\
\bottomrule
\end{tabular}
\end{center}

\begin{figure}[H]
\includegraphics{S4_F04_st02_spectrum_high}
\caption{The diurnal band after removing a one-week rolling mean. The daily cycle
and its second and third harmonics are clear.}
\end{figure}

The 7.28-day feature is a property of the analysis. A rolling mean of width $W$
acts as a high-pass whose response peaks near $W$ itself, and 168\,h is seven
days; re-running with a 720-hour window moves the feature, which is the test that
identifies it. A Welch periodogram on the longest unbroken block confirms the
daily peak independently at 1.0000\,d.

% ===========================================================================
\section{Question 2 --- decomposition, and whether it holds still}
\label{sec:harmonics}

\subsection{One limit that no amount of analysis removes}

Separating an independent \textbf{180.0-day} process from the annual second
harmonic at \textbf{182.62 days} requires resolving a frequency difference of
$8.0 \times 10^{-5}$ cycles per day. That needs a record of roughly
\textbf{34 years}. With six and a half it cannot be done, and no choice of
estimator, window or grid changes that.

Phase behaviour separates them where frequency cannot. If the semi-annual term is
harmonic distortion of the annual cycle --- which any nonlinear response to a
sinusoidal annual driver produces automatically --- its phase advances at exactly
twice the annual rate, so $2\varphi_{\mathrm{annual}} -
\varphi_{\mathrm{semiannual}}$ stays constant. An independent process drifts.

\subsection{Amplitudes, phases, and their drift}

\begin{center}
\small
\begin{tabular}{lR{1.8cm}R{1.9cm}R{1.6cm}R{2.4cm}R{1.6cm}}
\toprule
\textbf{Component} & \textbf{period} & \textbf{amplitude} & \textbf{s.e.} &
\textbf{peak, days after start} & \textbf{var.\ share} \\
\midrule
annual      & 365.25 d & \textbf{44.56} & 1.36 & 208.6 & 0.376 \\
semiannual  & 182.62 d & 13.99 & 1.34 & 112.8 & 0.037 \\
diurnal     &   1.00 d &  8.04 & 0.20 &   0.1 & 0.012 \\
semidiurnal &   0.50 d &  2.96 & 0.10 &   0.3 & 0.002 \\
\bottomrule
\multicolumn{6}{l}{\footnotesize Amplitudes in mdeg; Newey--West standard errors
with 24 lags. Model $R^2 = 0.487$.}
\end{tabular}
\end{center}

\begin{figure}[H]
\includegraphics{S4_F05_st02_harmonic_fit}
\caption{The harmonic model fitted across the gaps over the whole era. It passes
through the outages because a least-squares harmonic model does not require
continuity.}
\end{figure}

Refitting in two-year windows stepped by ninety days:

\begin{center}
\small
\begin{tabular}{lR{1.9cm}R{1.4cm}R{2.6cm}R{2.6cm}}
\toprule
\textbf{Component} & \textbf{mean amp.} & \textbf{s.d.} &
\textbf{amplitude range} & \textbf{phase range} \\
\midrule
annual     & 41.89 & 8.72 & 30.6 -- 53.3 & 189.5 -- 215.0 d \\
semiannual &  9.14 & 3.07 &  4.7 -- 13.9 &  87.7 -- 134.5 d \\
\bottomrule
\end{tabular}
\end{center}

\begin{figure}[H]
\includegraphics{S4_F06_st02_sliding_harmonics}
\caption{Amplitude and phase of the annual and semi-annual components along the
record. Neither is constant: the annual amplitude varies by a factor of 1.7 and
its phase by 25 days. This figure explains why the seasonal model fails to
predict in Section~\ref{sec:budget}.}
\end{figure}

\panel{The semi-annual term is the annual second harmonic}{%
The phase-lock test over twelve windows gives a circular standard deviation of
$2\varphi_{\mathrm{annual}} - \varphi_{\mathrm{semiannual}}$ of
\textbf{28.5\textdegree}, resultant length 0.884, against
\textbf{103.9\textdegree} for an unlocked null. The two components track each
other.

\medskip
The 183-day peak is therefore harmonic distortion of a non-sinusoidal annual
cycle, not a separate process with a six-month period. Its amplitude carries
information about the \emph{shape} of the annual response, not about any
independent semi-annual forcing.}

\begin{figure}[H]
\includegraphics{S4_F07_st02_diurnal}
\caption{Mean diurnal cycle over the whole era, for the target and both drivers.}
\end{figure}

% ===========================================================================
\section{Question 3 --- imputation}
\label{sec:imputation}

Filling a hole is defensible only up to the length at which the filler stops
knowing anything the data did not already contain. Forty randomly placed gaps per
length, eight lengths, four fillers.

\begin{center}
\small
\begin{tabular}{rR{2.0cm}R{2.4cm}R{1.8cm}R{2.2cm}}
\toprule
\textbf{gap} & \textbf{interpolate} & \textbf{seasonal naive} &
\textbf{harmonic} & \textbf{ridge, drivers} \\
\midrule
1 h   & \textbf{1.12} & 6.14 & 37.80 & 42.00 \\
3 h   & \textbf{1.59} & 4.51 & 28.34 & 32.32 \\
6 h   & \textbf{3.18} & 5.19 & 27.77 & 34.51 \\
12 h  & 6.07 & \textbf{5.37} & 25.44 & 34.99 \\
24 h  & 7.11 & \textbf{4.40} & 33.44 & 37.69 \\
72 h  & 8.15 & \textbf{6.90} & 27.78 & 37.03 \\
168 h & 8.73 & \textbf{8.27} & 25.71 & 32.74 \\
720 h & \textbf{8.23} & 22.93 & 22.44 & 32.17 \\
\bottomrule
\multicolumn{5}{l}{\footnotesize MAE in mdeg, against a signal standard deviation
of 51.6\,mdeg.}
\end{tabular}
\end{center}

\begin{figure}[H]
\includegraphics{S4_F08_st02_gaps_imputation}
\caption{Left: where the missing hours are, by gap length. Right: what each
filler costs against injected gaps of known length.}
\end{figure}

Applying a five-millidegree tolerance --- a tenth of the signal's standard
deviation --- gives \textbf{interpolation to 6\,h} and \textbf{seasonal-naive to
24\,h}; the model-based fillers never qualify.

Two results deserve emphasis. The \textbf{model-based fillers are the worst of
the four at every gap length}, by a wide margin: the harmonic model explains
under half the variance and knows nothing about local behaviour, and the
driver-based ridge is worse still --- which anticipates
Section~\ref{sec:budget}'s finding that these drivers carry almost no predictive
content. Simple interpolation beats both by a factor of eight at short gaps.

And the policy interacts with the gap inventory to give the practical answer:
filling everything up to 24 hours addresses gaps holding \textbf{9.6\,\%} of the
missing hours. The 84.5\,\% locked in stretches longer than thirty days is not
recoverable by any method tested, and should be treated as absent rather than
reconstructed.

% ===========================================================================
\section{Question 4 --- operators and the best single predictor}

\begin{figure}[H]
\includegraphics{S4_F09_st02_correlations}
\caption{Correlation structure of the legacy channels, on levels (left) and on
first differences (right). Levels correlations between two trending series are
inflated by the shared trend, so the differenced matrix is the honest one for a
signal with drift. Note that the raw and compensated inclination sit at opposite
signs against air temperature, which is Section~\ref{sec:sign} seen from a
second angle.}
\end{figure}

\subsection{Full band against diurnal band}

Before any driver is ranked it is given the chance to act through the operator
that suits it: a transport delay, and a thermal inertia. The delay is bounded at
twelve hours because both candidates are \textbf{external forcings}, which must
causally precede the response (\texttt{docs/raw-data-format.md} Section~7.5).

\begin{center}
\small
\begin{tabular}{lR{1.5cm}R{1.3cm}R{1.5cm}R{1.6cm}R{2.0cm}}
\toprule
\textbf{Driver} & \textbf{delay} & \textbf{$\tau$} & \textbf{$r$} &
\textbf{$R^2$} & \textbf{$R^2$ at $d=\tau=0$} \\
\midrule
\multicolumn{6}{l}{\itshape Full band} \\
\quad \texttt{tair} & 0 h &   0 h & $-0.629$ & 0.395 & 0.395 \\
\quad \texttt{rh}   & 0 h & 168 h & $+0.382$ & 0.146 & 0.091 \\
\addlinespace
\multicolumn{6}{l}{\itshape Diurnal band --- carried forward} \\
\quad \texttt{tair} & 0 h & 0 h & $-0.964$ & \textbf{0.929} & 0.929 \\
\quad \texttt{rh}   & 0 h & 0 h & $+0.595$ & 0.353 & 0.353 \\
\bottomrule
\end{tabular}
\end{center}

\begin{figure}[H]
\includegraphics{S4_F10_st02_operators_fullband}
\caption{Delay against inertia, full band. Relative humidity's optimum sits at
$\tau = 168$\,h, the last value in the grid.}
\end{figure}

\begin{figure}[H]
\includegraphics{S4_F11_st02_operators_diurnal}
\caption{The same grid on the diurnal band. Both optima collapse to zero delay
and zero time constant.}
\end{figure}

The comparison reproduces the extended-sensor study's finding exactly. On the
full band, relative humidity's optimum runs to the \emph{edge} of the $\tau$
grid, lifting explained variance from 0.091 to 0.146. A time constant of one week
is not a thermal property of a masonry wall; it is the width of filter needed to
turn a driver into a seasonal envelope, and it is pinned at the boundary because
a longer filter would fit better still. Band-limiting removes it. Reporting the
full-band operator as a physical time constant would have produced a confident
and spurious finding.

\subsection{The signed-delay control}

The extended-sensor study had to scan wall temperature over \emph{signed} delays,
because an internal state variable at unknown depth may legitimately lag the
deformation. This era has no wall probe: both candidates are external forcings,
so the signed scan here is a \textbf{control rather than a measurement}. If
either driver gained materially from a lead, the method would be at fault.

\begin{center}
\small
\begin{tabular}{lR{1.6cm}R{1.3cm}R{1.7cm}R{1.7cm}R{2.4cm}}
\toprule
\textbf{Driver} & \textbf{delay} & \textbf{$\tau$} & \textbf{$R^2$ signed} &
\textbf{$R^2$ causal} & \textbf{gained from lead} \\
\midrule
\texttt{tair} &  0 h & 0 h & 0.9291 & 0.9291 & 0.0000 \\
\texttt{rh}   & $-1$ h & 0 h & 0.3552 & 0.3534 & 0.0018 \\
\bottomrule
\end{tabular}
\end{center}

\begin{figure}[H]
\includegraphics{S4_F12_st02_operators_signed}
\caption{The diurnal-band grid over signed delays. Control only; feeds no model.}
\end{figure}

The control passes. Air temperature peaks exactly at zero. Relative humidity
peaks one hour early, but that lead buys \textbf{0.0018} of explained variance
--- an order of magnitude below the 0.01 threshold the study sets for a material
lead, and exactly what an optimum landing one grid step off zero on noise looks
like. Neither forcing leads the response in any meaningful sense, which is what
the physics requires.

\subsection{Ranking, on two windows}

Each driver is scored on its own, at its diurnal-band operator, by pooled error
over five rolling-origin folds. Autoregression is excluded: with it, everything
ties behind the autoregressive term and the ranking says nothing. Battery voltage
is carried as a \textbf{negative control}.

\begin{center}
\small
\begin{tabular}{lR{1.6cm}R{1.6cm}R{1.7cm}R{1.6cm}}
\toprule
\textbf{Driver} & \textbf{MAE} & \textbf{fold sd} & \textbf{$R^2$} &
\textbf{vs best} \\
\midrule
\multicolumn{5}{l}{\itshape Whole era} \\
\quad \texttt{tair}  & 36.70 & 11.19 & $+0.073$ & --- \\
\quad \texttt{batt}\,$^\dagger$ & 38.89 & 11.96 & $-0.215$ & $+6.0$\,\% \\
\quad \texttt{rh}    & 41.29 & 13.13 & $-0.327$ & $+12.5$\,\% \\
\addlinespace
\multicolumn{5}{l}{\itshape Unbroken block only} \\
\quad \texttt{tair}  & 38.04 & 16.19 & $-0.719$ & --- \\
\quad \texttt{rh}    & 41.83 & 23.51 & $-1.403$ & $+10.0$\,\% \\
\quad \texttt{batt}\,$^\dagger$ & 46.05 & 19.43 & $-1.583$ & $+21.1$\,\% \\
\bottomrule
\multicolumn{5}{l}{\footnotesize $^\dagger$ negative control.}
\end{tabular}
\end{center}

\begin{figure}[H]
\includegraphics{S4_F13_st02_rank}
\caption{Single-predictor skill on the whole era, without autoregression. Only
air temperature beats a constant, and only barely.}
\end{figure}

\panel{On the unbroken block, no driver beats a constant}{%
Air temperature's $R^2$ falls from $+0.073$ on the whole era to $\mathbf{-0.719}$
on the longest unbroken block --- worse than predicting the mean. The negative
control now ranks last, which is the one reassuring feature of the table.

\medskip
The marginal skill air temperature shows on the full era does not survive a
change of window. Whatever it was tracking was a property of the particular mix
of periods the whole record contains, not a stable relationship. This is the
clearest statement in the report that \textbf{the legacy drivers do not predict
this signal}, and Sections~\ref{sec:budget} and \ref{sec:sensitivity} reach it
independently.}

\panel{One discrepancy to read carefully}{%
Air temperature explains \textbf{92.9\,\%} of the variance of the \emph{diurnal
band} ($r = -0.964$) and \textbf{7.3\,\%} of the total variance out of sample.
Both numbers are correct and they are not in conflict: the diurnal band is a
small share of this signal's variance, which is dominated by seasonal and
lower-frequency structure an instantaneous driver cannot reach.

\medskip
Quoting the band-limited correlation as though it described the signal would
overstate what this channel knows by an order of magnitude.}

% ===========================================================================
\section{Question 5 --- combinations, and how to diagnose them}

\subsection{Collinearity}

\begin{center}
\small
\begin{tabular}{lR{1.6cm}R{2.6cm}l}
\toprule
\textbf{Channel} & \textbf{VIF} & \textbf{max $|r|$ with others} &
\textbf{closest} \\
\midrule
\texttt{tair} & 1.58 & 0.605 & \texttt{rh} \\
\texttt{rh}   & 1.58 & 0.605 & \texttt{tair} \\
\bottomrule
\end{tabular}
\end{center}

This is a useful contrast with the extended era, where air and wall temperature
were near-duplicates at $r = 0.924$ with variance inflation factors around eight.
Here the two candidates are only moderately related, so the collinearity caveat
that governed the sibling study's diagnostics does not bite: coefficients would
be readable. They are still not read, because the out-of-sample error is the
quantity that matters.

\subsection{Every combination}

With two candidates the exhaustive search is three subsets.

\begin{center}
\small
\begin{tabular}{lR{0.7cm}R{1.5cm}R{1.4cm}R{1.4cm}R{1.6cm}R{1.4cm}}
\toprule
\textbf{Subset} & \textbf{$k$} & \textbf{MAE} & \textbf{fold sd} &
\textbf{$R^2$} & \textbf{BIC} & \textbf{utility} \\
\midrule
\texttt{tair + rh} & 2 & \textbf{36.25} & 10.83 & $+0.095$ & 403\,197 & \textbf{0.685} \\
\texttt{tair}      & 1 & 36.70 & 11.19 & $+0.073$ & 403\,836 & 0.676 \\
\texttt{rh}        & 1 & 41.29 & 13.13 & $-0.327$ & 419\,943 & 0.601 \\
\bottomrule
\end{tabular}
\end{center}

\begin{figure}[H]
\includegraphics{S4_F14_st02_subsets}
\caption{Exhaustive subset search over the two legacy drivers.}
\end{figure}

The pair beats air temperature alone by \textbf{0.46\,mdeg, that is 1.3\,\%},
against a fold-to-fold spread of 10.8\,mdeg --- twenty-three times the
difference. \textbf{The ordering is not resolved by these data.} BIC and utility
agree with the error ranking, but all three are reading the same unresolved
margin.

\subsection{Which member is doing the work}

\begin{center}
\small
\begin{tabular}{lR{2.0cm}R{1.8cm}R{2.6cm}R{1.4cm}}
\toprule
\textbf{Channel} & \textbf{MAE without} & \textbf{gain} &
\textbf{permutation increase} & \textbf{sd} \\
\midrule
\texttt{tair} & 41.29 & 5.05 (12.2\,\%) & \textbf{9.07} & 2.69 \\
\texttt{rh}   & 36.70 & 0.46 (1.3\,\%)  & $-0.006$ & 0.124 \\
\bottomrule
\multicolumn{5}{l}{\footnotesize Subset \texttt{tair + rh}; full-subset MAE
36.25\,mdeg. All values in mdeg.}
\end{tabular}
\end{center}

\begin{figure}[H]
\includegraphics{S4_F15_st02_diagnostics}
\caption{Left: the cost of removing each channel and refitting. Right: the cost
of shuffling it with the model held fixed. Relative humidity sits at zero on the
right-hand panel, inside its own error bar.}
\end{figure}

The two diagnostics agree here, which is itself informative. Relative humidity's
leave-one-out gain is 0.46\,mdeg and its permutation increase is
$-0.006 \pm 0.124$ --- indistinguishable from zero. The fitted model does not
rely on it, and refitting without it costs nothing separable from fold noise.
Under the low VIF above there is no collinearity to hide behind: relative
humidity simply carries no content.

\subsection{The out-of-period test}
\label{sec:holdout}

Both models are fitted on the unbroken block and applied to the third block ---
2023-06-21 to 2024-05-04, on the far side of a 271-day outage --- without
refitting.

\begin{center}
\small
\begin{tabular}{lR{1.4cm}R{1.6cm}R{1.6cm}R{2.4cm}}
\toprule
\textbf{Model} & \textbf{$n$} & \textbf{MAE} & \textbf{RMSE} &
\textbf{vs cross-validated} \\
\midrule
\texttt{tair + rh} & 6997 & \textbf{62.49} & 65.00 & $\times 1.72$ \\
\texttt{tair} alone & 6997 & 63.56 & 65.92 & $\times 1.73$ \\
\bottomrule
\end{tabular}
\end{center}

\panel{Neither model transfers}{%
Both roughly \textbf{double} their cross-validated error, from 36 to 62--64\,mdeg
--- worse than the 48.8\,mdeg a constant achieves in the error budget of
Section~\ref{sec:budget}. Across a 271-day outage, a model fitted on
environmental drivers does not merely degrade; it becomes worse than predicting
the mean.

\medskip
The pair is marginally better than the single predictor here, reversing the
extended-sensor study's holdout result --- but by 1.07\,mdeg on two models that
are both failing, which settles nothing.}

\subsection{The diagnostic protocol}

The procedure used above is the answer to ``how is the best combination
diagnosed''. It is the sibling study's protocol with one rule added --- rule 2 ---
which this era's spectral work forced.

\begin{center}
\small
\begin{tabularx}{\textwidth}{R{0.5cm}L{5.6cm}Y}
\toprule
& \textbf{Rule} & \textbf{Failure it prevents} \\
\midrule
1 & Bound the operator scan by the mechanism and by half the dominant period &
A physically impossible optimum, or an aliased lead read as a delay \\
2 & Estimate the spectrum before imputing anything, and control for the sampling
pattern &
Letting the filling method decide which frequencies exist; reading an outage
pattern as a structural cycle \\
3 & Fit each driver operator on the band of interest &
A boundary optimum reported as a physical time constant \\
4 & Score without autoregression &
Every driver tying behind the autoregressive term \\
5 & Rolling-origin folds, report the spread &
A ranking that depends on where the single split fell \\
6 & Exhaustive subset search when the candidate set is small &
Selection bias and instability under collinearity \\
7 & Gain and permutation importance, never coefficients under high VIF &
Importance attributed by collinearity rather than content \\
8 & Weight by availability over the whole record &
Choosing a channel absent when it is needed \\
9 & Carry a negative control &
Mistaking shared diurnal shape for predictive content \\
10 & Confirm on an out-of-period block &
A combination that fits one period and transfers to none \\
\bottomrule
\end{tabularx}
\end{center}

On this record rules 5, 9 and 10 did the work: the fold spread showed the subset
ordering to be unresolved, the negative control ranked above relative humidity on
the whole era, and the out-of-period test showed both models failing.

% ===========================================================================
\section{Question 6 --- horizon, uncertainty and the error budget}
\label{sec:budget}

\subsection{The ladder}

The ladder adds one component at a time and scores every rung at every horizon,
so each component's contribution is read off consecutive rows. Every rung is
fitted before a chronological cut and scored across the whole test region at an
exact lead time, pooled over three origins.

\begin{center}
\small
\begin{tabular}{lR{1.2cm}R{1.2cm}R{1.2cm}R{1.2cm}R{1.2cm}R{1.3cm}R{1.3cm}}
\toprule
\textbf{Rung} & \textbf{1 h} & \textbf{24 h} & \textbf{168 h} &
\textbf{720 h} & \textbf{2160 h} & \textbf{4380 h} & \textbf{8760 h} \\
\midrule
climatology  & 48.80 & 48.82 & 48.98 & 49.56 & 50.57 & 51.07 & \textbf{51.67} \\
+ annual     & 50.12 & 50.15 & 50.32 & 50.96 & 52.42 & 55.12 & 55.43 \\
+ semiannual & 49.44 & 49.46 & 49.56 & 50.01 & 51.58 & 54.75 & 54.98 \\
+ diurnal    & 49.34 & 49.36 & 49.46 & 49.90 & 51.47 & 54.70 & 54.98 \\
+ AR         & \textbf{1.60} & \textbf{4.90} & \textbf{8.35} &
\textbf{12.27} & \textbf{23.49} & \textbf{45.03} & 61.55 \\
+ drivers    & 1.60 & 4.95 & 8.13 & 12.61 & 25.19 & 46.96 & 62.87 \\
\bottomrule
\multicolumn{8}{l}{\footnotesize MAE in mdeg. Horizons 3, 6, 12 and 72 h omitted
for width; they interpolate smoothly.}
\end{tabular}
\end{center}

\begin{figure}[H]
\includegraphics{S4_F16_st02_error_budget}
\caption{Left: error against horizon per rung. The calendar-only rungs are flat,
because a function of time does not degrade with lead; the autoregressive rungs
climb. Right: the error reduction attributed to each component.}
\end{figure}

\textbf{Autoregression is almost the entire budget.} At one hour it removes 47.74
of the 48.80\,mdeg climatological error, and it remains dominant out to 720\,h ---
one month --- where it still removes 37.63. It continues to beat climatology at
4380\,h (45.03 against 51.07) and finally fails at one year.

\textbf{The seasonal terms buy nothing, and the annual term is actively
harmful.} Adding it makes the forecast \emph{worse} by 1.33\,mdeg at every
horizon; the semi-annual term recovers 0.68 and the diurnal term 0.10. The reason
is Figure~6: the annual amplitude ranges over a factor of 1.7 and its phase by 25
days, so a fixed harmonic fitted on the training block does not describe the test
block. The cycle is real --- the spectrum is unambiguous --- but it is not
stationary, and only a stationary component can be extrapolated.

\textbf{The drivers add nothing measurable}, at every horizon, within
$\pm$0.25\,mdeg.

\subsection{What the trend term costs}

\begin{center}
\small
\begin{tabular}{lR{1.6cm}R{1.6cm}R{1.6cm}R{1.6cm}R{1.6cm}}
\toprule
\textbf{Model} & \textbf{1 h} & \textbf{24 h} & \textbf{720 h} &
\textbf{4380 h} & \textbf{8760 h} \\
\midrule
seasonal, no trend & 49.34 & 49.36 & 49.90 & 54.70 & 54.98 \\
seasonal + trend   & 134.34 & 134.38 & 135.65 & 143.92 & 148.15 \\
\midrule
\textbf{cost}      & \textbf{+85.00} & \textbf{+85.02} & \textbf{+85.75} &
\textbf{+89.22} & \textbf{+93.18} \\
\bottomrule
\end{tabular}
\end{center}

\panel{Do not extrapolate a fitted linear trend on this signal}{%
A straight line fitted on part of this record and extended across the rest costs
between 85 and 93\,mdeg at every horizon --- against a signal standard deviation
of 51.6\,mdeg and a climatological error of 49. It is by a wide margin the most
damaging single term available, and it is a term the grey-box pipeline includes
by default.

\medskip
The cause is the same non-stationarity that defeats the annual harmonic. A drift
term is defensible as a \emph{description} of a fitted window; as an
\emph{extrapolation} it is not.}

\subsection{Knowing the drivers in advance}

\begin{center}
\small
\begin{tabular}{rR{1.6cm}R{1.6cm}R{1.6cm}}
\toprule
\textbf{horizon} & \textbf{lagged} & \textbf{future} & \textbf{gain} \\
\midrule
1 h    &  1.60 &  1.61 & $-0.5$\,\% \\
24 h   &  4.95 &  4.26 & $+14.0$\,\% \\
168 h  &  8.13 &  4.99 & $+38.6$\,\% \\
720 h  & 12.61 &  9.43 & $+25.2$\,\% \\
8760 h & 62.87 & 61.82 & $+1.7$\,\% \\
\bottomrule
\end{tabular}
\end{center}

Perfect foreknowledge of the drivers is worth nothing at one hour, most at one
week, and almost nothing again at one year. At short horizons the signal's own
history already contains what the drivers would say; at long horizons neither
knows anything.

\subsection{NeuralProphet at the short horizons}
\label{sec:np}

NeuralProphet cannot span a year, so it is run on the longest unbroken block at
horizons to one week.

\panel{Correction: the annual term was never enabled}{%
An earlier version of this section stated that this run carried the annual term,
and qualified the result by noting that the block holds only 1.69 cycles so the
component would be weakly identified. That qualification described something
that does not exist. The function producing these numbers,
\texttt{ip.run\_horizon\_experiment}, fixed \texttt{yearly\_seasonality} to
false, so the configuration below carries a daily seasonality, an autoregressive
term and the drivers, and no annual component whatever.

\medskip
The numbers in this section are unaffected --- they always described the model
that was actually fitted. What was wrong was the description. The parameter is
now explicit, and Section~\ref{sec:usable} enables it on a span long enough to
identify it.}

\begin{center}
\small
\begin{tabular}{rR{1.4cm}R{1.6cm}R{1.9cm}R{1.5cm}R{1.6cm}}
\toprule
\textbf{$h$} & \textbf{MAE} & \textbf{persist.} & \textbf{skill vs p.} &
\textbf{cover.} & \textbf{width} \\
\midrule
\multicolumn{6}{l}{\itshape lagged --- the deployable configuration} \\
1   &  3.35 & 1.55 & $-1.160$ & 100.0\,\% & 164.2 \\
12  &  4.76 & 6.74 & $+0.294$ & 100.0\,\% & 165.5 \\
24  &  5.66 & 4.41 & $-0.283$ & 100.0\,\% & 167.0 \\
168 & 13.53 & 7.59 & $-0.783$ &  99.4\,\% & 185.3 \\
\addlinespace
\multicolumn{6}{l}{\itshape future --- regressors known over the horizon} \\
1   & 2.01 & 1.54 & $-0.305$ & 61.5\,\% & 14.5 \\
12  & 3.15 & 6.65 & $+0.526$ & 60.6\,\% & 15.7 \\
24  & 3.56 & 4.38 & $+0.188$ & 61.2\,\% & 16.8 \\
168 & 6.27 & 7.43 & $+0.156$ & 65.0\,\% & 21.9 \\
\bottomrule
\end{tabular}
\end{center}

\begin{figure}[H]
\includegraphics{S4_F17_st02_np_skill}
\caption{NeuralProphet error and skill against persistence on the longest
unbroken block.}
\end{figure}

The deployable configuration beats persistence only between 6 and 12 hours, and
beats the seasonal-naive baseline only out to 6 hours. With perfect future
regressors the useful band extends to a full week.

\begin{figure}[H]
\includegraphics{S4_F19_st02_fan}
\caption{Ten days of the test period at a twenty-four-hour horizon, with the
nominal 90\,\% interval. The interval is so wide that the observed series never
approaches its edges.}
\end{figure}

\subsection{Is the quoted uncertainty honest?}

\begin{figure}[H]
\includegraphics{S4_F18_st02_np_uncertainty}
\caption{Calibration and width of the nominal 90\,\% interval. The lagged
configuration covers 100\,\% of observations --- by producing intervals more than
three times the signal's own standard deviation.}
\end{figure}

\panel{Both configurations fail calibration, in opposite directions}{%
The \textbf{lagged} model covers \textbf{100\,\%} of observations against a
nominal 90, which is not a success: it achieves this with a median interval width
of 164--185\,mdeg, more than three times the signal's standard deviation of 51.6.
An interval that wide excludes nothing and so detects nothing.

\medskip
The \textbf{future} model covers \textbf{61--65\,\%} with widths of 14--22\,mdeg
--- informative in width, but substantially under-covering.

\medskip
The extended-sensor study found its deployable configuration under-covering at
54--67\,\%. Here the same configuration over-covers by inflating the interval.
Both are miscalibrated, in opposite directions, which is a reason to recalibrate
by a split-conformal procedure rather than to trust either quantile head.}

\subsection{What the joint multi-horizon head costs}

One fit produced every horizon. Refitting with \texttt{n\_forecasts} set to
exactly one horizon isolates the cost of sharing weights across 168 output steps.

\begin{center}
\small
\begin{tabular}{R{1.6cm}R{2.2cm}R{2.2cm}R{2.4cm}}
\toprule
\textbf{$h$ [h]} & \textbf{dedicated} & \textbf{joint head} &
\textbf{joint penalty} \\
\midrule
1  & 2.93 & 3.35 & $+14.3$\,\% \\
24 & 5.96 & 5.66 & $-5.0$\,\% \\
\bottomrule
\end{tabular}
\end{center}

The joint head costs 14\,\% at one hour and \emph{gains} 5\,\% at twenty-four ---
the same pattern the extended-sensor study found, consistent with shared weights
acting as a regulariser at longer horizons and a distraction at the shortest. The
conclusion is unchanged: even the dedicated one-hour model, at 2.93\,mdeg, does
not beat persistence at 1.55.

% ===========================================================================
\section{Sensitivity}
\label{sec:sensitivity}

The headline is the autoregressive rung's error one day ahead, 4.90\,mdeg.

\begin{center}
\small
\begin{tabular}{lR{1.6cm}R{1.6cm}R{1.8cm}}
\toprule
\textbf{Variation} & \textbf{value} & \textbf{change} & \textbf{change \%} \\
\midrule
baseline                         &  4.90 & --- & --- \\
compensation coefficient 0.010   & 13.88 & $+8.98$ & \textbf{$+183.2$\,\%} \\
model block only                 &  9.46 & $+4.56$ & $+93.0$\,\% \\
compensation coefficient 0.00163 &  2.07 & $-2.83$ & \textbf{$-57.7$\,\%} \\
autoregressive depth 48          &  4.78 & $-0.12$ & $-2.4$\,\% \\
autoregressive depth 12          &  4.98 & $+0.08$ & $+1.6$\,\% \\
origins 0.4 / 0.5 / 0.6          &  4.92 & $+0.02$ & $+0.4$\,\% \\
no seasonal terms at all         &  4.89 & $-0.01$ & $-0.2$\,\% \\
no semi-annual term              &  4.90 & $-0.00$ & $-0.1$\,\% \\
drivers dropped                  &  4.90 & $ 0.00$ & \textbf{$0.0$\,\%} \\
\bottomrule
\end{tabular}
\end{center}

\begin{figure}[H]
\includegraphics{S4_F20_st02_sensitivity}
\caption{Sensitivity of the 24-hour autoregressive error to the study's choices.}
\end{figure}

\textbf{The compensation coefficient is the most consequential decision in the
study}, moving the headline by $-58$\,\% to $+183$\,\%. Every absolute error in
this report is a statement about a target built with a particular coefficient,
and that coefficient has never been calibrated for this instrument.

\textbf{The seasonal decomposition contributes nothing to prediction.} Removing
every seasonal term changes the 24-hour error by 0.2\,\%. Sections 3 and 4
establish that the cycles are real, identified and interpretable; this table and
Section~\ref{sec:budget} establish that they are useless for forecasting, because
they are not stationary. Both statements are true and the report keeps them
separate.

\textbf{The drivers are exactly worthless.} Dropping them changes the headline by
0.00\,\% --- the third independent route to that conclusion, after the
unbroken-block ranking and the out-of-period test.

% ===========================================================================
\section{How much of the record can actually be used}
\label{sec:usable}

Section~\ref{sec:record} divided the study on a premise: that the spectral and
harmonic work tolerates gaps by construction while the autoregressive work does
not, and must therefore run on the longest unbroken block. The first half is
right. The second half is too strong, and every result from
Section~\ref{sec:budget} onwards was computed under it.

An autoregressive model does not need a continuous series. It needs windows of
$n_{\text{lags}}$ consecutive observations followed by $n_{\text{forecasts}}$
more, and those windows exist in every block, not only in the longest one. The
restriction was never tested; this section tests it.

That also reframes what imputation is for. Filling a hole is not a way of making
a record continuous --- it is a way of buying back the windows the hole
destroyed. A single missing hour costs a full window's worth of samples however
brief it is, so the many short interruptions are cheap to fill and expensive to
leave, and a months-long outage is the reverse.

\subsection{What each gap tolerance buys}

Section~\ref{sec:record} swept the tolerance to three days. Extended, it shows
where the record actually divides.

\begin{center}
\small
\begin{tabular}{rR{1.4cm}R{1.6cm}R{1.4cm}R{2.1cm}l}
\toprule
\textbf{tolerance} & \textbf{blocks} & \textbf{longest} &
\textbf{cycles} & \textbf{fabricated} & \textbf{two cycles?} \\
\midrule
6 h    & 11 &  211.5 d & 0.58 &  1.6\,\% & \fail \\
24 h   &  6 &  600.2 d & 1.64 &  2.6\,\% & \fail \\
72 h   &  6 &  617.2 d & 1.69 &  3.2\,\% & \fail \\
7 d    &  6 &  617.2 d & 1.69 &  3.2\,\% & \fail \\
30 d   &  5 &  617.2 d & 1.69 &  3.2\,\% & \fail \\
45 d   &  3 &  617.2 d & 1.69 &  3.2\,\% & \fail \\
\textbf{60 d} & 2 & \textbf{1353.9 d} & \textbf{3.71} &
\textbf{14.3\,\%} & \pass \\
90 d   &  2 & 1353.9 d & 3.71 & 14.3\,\% & \pass \\
\bottomrule
\end{tabular}
\end{center}

Between three days and forty-five the longest block does not move at all. The
record holds exactly seven interruptions longer than three days, and their
positions explain the whole table:

\begin{center}
\small
\begin{tabular}{llR{1.6cm}}
\toprule
\textbf{From} & \textbf{To} & \textbf{Length} \\
\midrule
2018-09-23 & 2018-11-21 & 59.0 d \\
2019-01-14 & 2019-02-04 & 20.8 d \\
2019-04-20 & 2019-05-11 & 20.5 d \\
2020-06-04 & 2020-07-31 & 57.3 d \\
\textbf{2022-04-09} & \textbf{2023-06-21} & \textbf{437.6 d} \\
2024-05-04 & 2024-06-14 & 41.5 d \\
2024-08-27 & 2024-10-09 & 42.7 d \\
\bottomrule
\end{tabular}
\end{center}

The four in the earlier half stand between the record's start and the 617-day
block, so absorbing any subset of them short of all four leaves that block
unchanged --- which is why the longest span is flat from three days to
forty-five. The binding constraint is the \textbf{59-day} interruption of autumn
2018, and the table jumps precisely when the tolerance passes it: at
\textbf{sixty days} the whole earlier half consolidates into one stretch of
1353.9 days --- \textbf{3.71 annual cycles}, the first window in this project to
clear the two-cycle rule --- at the price of inventing 14.3\,\% of it.

\textbf{The record is permanently in two halves.} The 438-day outage between
April 2022 and June 2023 never closes at any tolerance a person would defend. The
earlier half reaches 3.71 cycles; the later half reaches 1.67 and never
qualifies.

\subsection{Windows gained against hours invented}

Days are not what a model consumes. Counting windows instead, at the
configuration Section~\ref{sec:budget} actually used
($n_{\text{lags}} = 24$, $n_{\text{forecasts}} = 168$):

\begin{center}
\small
\begin{tabular}{lR{1.3cm}R{1.5cm}R{1.5cm}R{1.8cm}R{1.5cm}}
\toprule
\textbf{Policy} & \textbf{cycles} & \textbf{windows} & \textbf{vs block} &
\textbf{invented} & \textbf{per hour} \\
\midrule
model block (control)  & 1.69 & 14\,641 & 1.00 &   492 h & --- \\
full era, no fill      & 6.57 & 17\,128 & 1.17 &     0 h & --- \\
\textbf{full era, fill 6 h} & \textbf{6.57} & \textbf{33\,074} &
\textbf{2.26} & \textbf{881 h} & \textbf{20.9} \\
full era, fill 24 h    & 6.57 & 39\,345 & 2.69 & 1\,756 h & 14.1 \\
full era, fill 72 h    & 6.57 & 39\,792 & 2.72 & 1\,821 h & 13.8 \\
\bottomrule
\end{tabular}
\end{center}

\begin{figure}[H]
\includegraphics{S4_F21_st02_window_accounting}
\caption{Left: usable autoregressive windows under each policy, annotated with
the annual cycles the span contains. Right: the share of that span which had to
be invented to obtain them.}
\end{figure}

\panel{Filling the short interruptions is the best-value operation in the study}{%
Taking the whole era \emph{without filling anything} buys only 17\,\% more
windows than the block, because 554 separate interruptions each destroy a full
window. Filling just the 475 holes of six hours or less --- 881 hours, or
\textbf{1.5\,\% of the span}, at a measured interpolation error of 1.1 to
3.2\,mdeg against a signal standard deviation of 51.6 --- raises that to
\textbf{2.26 times} the block's windows, and buys \textbf{20.9 windows for every
hour invented}.

\medskip
Past six hours the return falls and never recovers: the 24-hour policy buys 14.1
windows per invented hour and the 72-hour policy 13.8. The six-hour limit is also
exactly what the Section~\ref{sec:imputation} benchmark licenses for
interpolation, so the efficient choice and the defensible one coincide.}

\subsection{Is this one series, or two?}

A 438-day outage is long enough for an instrument to be serviced, remounted or
replaced, and an inclinometer's absolute level means nothing across such an
event. Fitted straight through it, a trend term would read any offset as
structural movement.

The naive comparison confounds the offset with the season, since the last
observations before the outage and the first after it sit at opposite ends of the
annual cycle.

\begin{center}
\small
\begin{tabular}{lR{2.0cm}R{2.0cm}R{1.8cm}R{1.4cm}}
\toprule
\textbf{Test} & \textbf{before} & \textbf{after} & \textbf{step} &
\textbf{$t$} \\
\midrule
response, raw        & 151.66 &  91.55 & $-60.11$ & $-71.6$ \\
response, deseasoned &   9.28 &  22.70 & $+13.42$ & $+22.9$ \\
control \texttt{tair} &  7.98 &  24.58 & $+16.61$ & $+63.5$ \\
\bottomrule
\end{tabular}
\end{center}

\noindent{\footnotesize Response in mdeg, control in \degC. Thirty days each
side; the deseasoned row has the Section~\ref{sec:harmonics} harmonic fit
subtracted.}

The raw step of $-60$\,mdeg is mostly calendar: the control is 16.6\,\degC\
warmer after the outage than before it, and the compensated inclination moves
against temperature. Once the seasonal expectation is removed, \textbf{13.4\,mdeg
of step survives} --- about a quarter of the signal's standard deviation, and far
too large to dismiss.

\textbf{So the era is treated as two series, not one.} That is not a compromise:
it is what the prediction regime below does anyway, because it presents each
complete stretch to the model separately and normalises each one on its own. An
offset introduced during an outage therefore cannot propagate into the fit at
all.

\subsection{Decomposition: does a longer window help?}

The companion manuscript states as a framework rule that a window of one to two
annual cycles suffices to fit a yearly Fourier component reliably. This record is
the only place in the project able to test it. The same NeuralProphet
configuration is fitted on 1.69 cycles and on 6.57, and both are measured against
the harmonic fit of Section~\ref{sec:harmonics}, whose single fixed basis gives
an annual signal exactly one place to go.

\begin{center}
\small
\begin{tabular}{lR{2.4cm}R{2.4cm}}
\toprule
& \textbf{model block} & \textbf{full era} \\
& \textbf{1.69 cycles} & \textbf{6.57 cycles} \\
\midrule
observations                        & 14\,340 & 39\,499 \\
\midrule
sd of trend [mdeg]                  & 23.29 & 35.05 \\
sd of yearly term [mdeg]            & 35.07 & 30.07 \\
\textbf{trend's share}              & \textbf{39.9\,\%} & \textbf{53.8\,\%} \\
\midrule
annual amplitude [mdeg]             & 48.37 & 45.48 \\
harmonic reference [mdeg]           & 42.54 & 43.31 \\
\textbf{amplitude ratio}            & \textbf{1.137} & \textbf{1.050} \\
phase gap to reference [d]          & 10 & 29 \\
\bottomrule
\end{tabular}
\end{center}

\begin{figure}[H]
\includegraphics{S4_F22_st02_decomposition}
\caption{The whole-era decomposition. Top: the observations with the harmonic fit
over them; below, the trend and each seasonal component NeuralProphet separated.}
\end{figure}

\panel{More cycles fix the amplitude and make the trend worse}{%
The longer span cuts the annual amplitude's overestimate from \textbf{13.7\,\% to
5.0\,\%}. On that measure the framework rule is too optimistic: 1.69 cycles
identified an annual term nearly fourteen per cent too large, and the extra four
cycles largely repaired it.

\medskip
On every other measure the longer span is no better. The trend's share of the
low-frequency variance rises from \textbf{39.9\,\% to 53.8\,\%} --- over the
whole era NeuralProphet's piecewise-linear trend holds \emph{more} variance than
the annual cycle it is supposed to sit beneath --- and the phase drifts from
10 days off the harmonic reference to 29.

\medskip
The finding is therefore not "use a longer window" but something sharper:
\textbf{NeuralProphet's trend is not a safe container at either length}. At both
spans it absorbs a large share of an annual signal that a fixed harmonic basis
attributes correctly, which is the same failure the
\texttt{thermal\_compensation\_legacy} study recorded on air temperature and the
same term Section~\ref{sec:budget} priced at 85--93\,mdeg when extrapolated. A
decomposition that must apportion trend against season correctly should not be
read off this component alone.}

\subsection{Prediction across every complete stretch}

\panel{NeuralProphet 0.8.0's \texttt{drop\_missing} does not work}{%
The library's documented mechanism for exactly this situation is
\texttt{drop\_missing}, which is meant to discard the autoregressive windows
touching a gap and train on the rest. In version 0.8.0 it does neither correctly.
Training on a frame that actually requires a drop returns a \textbf{NaN loss from
the first epoch} and leaves every weight NaN; the prediction path raises
\texttt{Length of values ... does not match length of index}, because the rows it
drops are not dropped from the frame the predictions are written into. Both were
verified on this record.

\medskip
The route used instead is the library's multi-series support, which is sound: the
surviving stretches are presented through the \texttt{ID} column, one set of
shared components is fitted across all of them, and normalisation is per series.
The seasonal terms therefore see all six and a half years while no autoregressive
window is ever formed across an outage.}

\begin{center}
\small
\begin{tabular}{rR{1.3cm}R{1.4cm}R{1.3cm}R{1.9cm}R{1.9cm}}
\toprule
\textbf{$h$} & \textbf{block} & \textbf{full era} & \textbf{change} &
\textbf{coverage} & \textbf{width} \\
\midrule
\multicolumn{6}{l}{\itshape lagged --- the deployable configuration} \\
1   &  3.35 &  1.88 & $-43.8$\,\% & 100 $\to$ 84 & 164 $\to$ 12 \\
3   &  4.55 &  3.19 & $-29.7$\,\% & 100 $\to$ 80 & 170 $\to$ 14 \\
6   &  4.41 &  4.15 & $-5.9$\,\%  & 100 $\to$ 79 & 167 $\to$ 16 \\
12  &  4.76 &  5.00 & $+5.1$\,\%  & 100 $\to$ 77 & 165 $\to$ 17 \\
24  &  5.66 &  5.53 & $-2.3$\,\%  & 100 $\to$ 76 & 167 $\to$ 19 \\
72  &  8.92 &  8.69 & $-2.5$\,\%  & 100 $\to$ 65 & 175 $\to$ 23 \\
168 & 13.53 & 11.12 & $-17.8$\,\% &  99 $\to$ 55 & 185 $\to$ 23 \\
\addlinespace
\multicolumn{6}{l}{\itshape future --- regressors known over the horizon} \\
1   & 2.01 & 1.08 & $-46.1$\,\% & 62 $\to$ 82 & 14 $\to$ 10 \\
3   & 2.31 & 1.72 & $-25.6$\,\% & 61 $\to$ 79 & 15 $\to$ 10 \\
6   & 2.79 & 1.86 & $-33.5$\,\% & 61 $\to$ 79 & 16 $\to$ 10 \\
12  & 3.15 & 2.04 & $-35.1$\,\% & 61 $\to$ 78 & 16 $\to$ 10 \\
24  & 3.56 & 2.11 & $-40.7$\,\% & 61 $\to$ 77 & 17 $\to$ 10 \\
72  & 4.56 & 2.84 & $-37.7$\,\% & 63 $\to$ 76 & 20 $\to$ 12 \\
168 & 6.27 & 3.51 & $-44.0$\,\% & 65 $\to$ 76 & 22 $\to$ 13 \\
\bottomrule
\end{tabular}
\end{center}

\noindent{\footnotesize MAE in mdeg; coverage in per cent against a nominal 90;
width in mdeg. Each pair reads block $\to$ full era.}

\begin{figure}[H]
\includegraphics{S4_F23_st02_segment_skill}
\caption{NeuralProphet fitted across every complete stretch of the era, against
the persistence and seasonal-naive baselines.}
\end{figure}

\panel{The uncertainty was an artefact of the window, not of the model}{%
Section~\ref{sec:np} reported that the deployable configuration covered
\textbf{100\,\%} of observations with intervals of \textbf{164--185\,mdeg} ---
more than three times the signal's own standard deviation, so wide that they
excluded nothing and detected nothing. That was the single worst result in this
report.

\medskip
Fitted across the whole record instead, the same configuration produces intervals
of \textbf{12--23\,mdeg} at a coverage of \textbf{55--84\,\%}. The intervals are
now narrower than the signal and are no longer trivially satisfied; they
under-cover at the long horizons, which a split-conformal recalibration can
correct and an interval three times too wide cannot.

\medskip
The point estimate improves too, by 44\,\% one hour ahead and 18\,\% one week
ahead, and the useful band against persistence widens from 6--12 h to 1--12 h.
With known future regressors every horizon improves, by 26 to 46\,\%.

\medskip
None of this comes from a better model. It is the same architecture, the same
depth, the same drivers and the same number of epochs. It comes from 2.26 times
the training windows over 3.9 times the seasonal span, bought with
\textbf{1.5\,\% fabricated data}.}

\subsection{The bridged block, and what the fill costs it}

A genuinely continuous series of two annual cycles can be had only by inventing
the interruptions inside it. The block is built here because the question was
asked, and it is reported with its cost attached: 1353.9 days at 3.71 cycles, of
which \textbf{14.3\,\% is fabricated}, including six holes longer than anything
the Section~\ref{sec:imputation} benchmark licenses and one of 59 days.

Fitting the harmonic model twice --- once on the bridged series, once with the
fabricated positions masked out again --- prices the distortion directly.

\begin{center}
\small
\begin{tabular}{lR{2.0cm}R{2.4cm}R{1.8cm}}
\toprule
\textbf{Component} & \textbf{bridged} & \textbf{observed only} &
\textbf{bias} \\
\midrule
annual       & 43.41 & 44.93 & $-3.4$\,\% \\
semi-annual  & 13.17 & 15.86 & $\mathbf{-17.0}$\,\% \\
diurnal      &  7.16 &  8.43 & $\mathbf{-15.0}$\,\% \\
semi-diurnal &  2.49 &  2.91 & $\mathbf{-14.6}$\,\% \\
\bottomrule
\end{tabular}
\end{center}

\noindent{\footnotesize Amplitudes in mdeg.}

\textbf{The fill suppresses every component, and the short ones most.} A straight
line drawn across a 59-day hole carries the annual level tolerably --- the annual
amplitude falls by only 3.4\,\% --- but it contains no diurnal variation at all,
so the components with periods shorter than the hole lose 15 to 17\,\% of their
amplitude. The bias is systematically downward, because interpolation adds span
without adding oscillation.

\textbf{This block is not recommended and nothing else here depends on it.}
Regime~B obtains 6.57 annual cycles for 1.5\,\% fabricated data; this obtains
3.71 for 14.3\,\%. The bridged block would be worth its cost only if a strictly
uninterrupted series were required for some other reason, and no analysis in this
report requires one.

% ===========================================================================
\section{Verdicts}

\begin{center}
\small
\begin{tabularx}{\textwidth}{L{4.0cm}Y}
\toprule
\textbf{Question} & \textbf{Answer} \\
\midrule
1 · periodicities &
Annual at 370.1\,d (power 0.400), semi-annual at 183.0\,d, diurnal at 23.93\,h
with second and third harmonics, and an unanticipated component at 113.9\,d. All
clear a 200-replicate permutation null. \\
1 · spurious peaks &
The 971-day feature has more window-function power than signal power: an outage
artefact. The 7.28-day feature is an artefact of the detrending filter. \\
2 · is 180 d independent? &
\textbf{No.} Phase-locked to the annual term at 28.5\textdegree\ against
103.9\textdegree\ for an unlocked null: the annual second harmonic. Frequency
alone could never separate them --- that needs 34 years. \\
2 · best decomposition window &
The whole era, necessarily. No contiguous block holds two annual cycles; the
longest holds 1.69. \\
2 · are the components stable? &
\fail. Annual amplitude 30.6--53.3\,mdeg, phase spanning 25 days. \\
3 · imputation limit &
Interpolation to 6\,h, seasonal-naive to 24\,h, model-based fillers never. This
addresses 9.6\,\% of the missing hours; the 84.5\,\% in gaps over 30 days is not
recoverable. \\
4 · best single predictor &
\texttt{tair}, at $R^2 = +0.073$ on the whole era. \fail\ on the unbroken block,
where no driver beats a constant ($R^2 = -0.719$). \\
4 · lag structure &
Both forcings instantaneous. The signed control passes: the only negative
optimum is \texttt{rh} at $-1$\,h, buying 0.0018 $R^2$. \\
4 · negative control &
\texttt{batt} ranks \emph{second of three} on the whole era, above relative
humidity --- the noise floor sits above one of the two real channels. \\
5 · best combination &
\texttt{tair + rh}, by 1.3\,\% against a fold spread twenty-three times larger.
\fail\ as a resolved ordering. \\
5 · out-of-period &
\fail. Both models roughly double their error, to 62--64\,mdeg --- worse than the
48.8\,mdeg a constant achieves. \\
6 · error budget &
Autoregression is 47.7 of the 48.8\,mdeg at 1\,h and still dominates at one
month. Seasonal terms contribute 0.1\,mdeg; drivers contribute nothing. \\
6 · useful horizon &
Autoregression beats climatology out to 4380\,h (six months). NeuralProphet's
deployable configuration beats persistence only from 6 to 12\,h. \\
6 · the trend term &
\fail. Extrapolating it costs 85--93\,mdeg at every horizon. \\
6 · uncertainty &
\fail\ both ways \emph{on the unbroken block}: the deployable configuration
covers 100\,\% with intervals over three times the signal's standard deviation;
the perfect-foresight one covers 61--65\,\%. Section~\ref{sec:usable} largely
repairs this. \\
\bottomrule
\end{tabularx}
\end{center}

\noindent Question 7 was raised by the answers above and is
Section~\ref{sec:usable}.

\begin{center}
\small
\begin{tabularx}{\textwidth}{L{4.0cm}Y}
\toprule
\textbf{Question} & \textbf{Answer} \\
\midrule
7 · must autoregression use one block? &
\textbf{No}, and the assumption was expensive. A model needs windows, not a
continuous series, and they exist in every block. \\
7 · usable record &
The longest block gives 1.69 annual cycles and 14\,641 windows. Fitting across
every complete stretch gives \textbf{6.57 cycles and 33\,074 windows} ---
\textbf{2.26 times} as many --- for \textbf{1.5\,\%} fabricated data. \\
7 · where imputation pays &
Only below six hours: 475 short gaps cost 881 h and buy \textbf{21 windows per
invented hour}. At 24 h the return falls to 14.1, and the six-hour limit is also
what the benchmark licenses. \\
7 · one series or two? &
\textbf{Two.} The 438-day outage shows a $-60.1$\,mdeg step, of which
$+13.4$\,mdeg survives removing the seasonal expectation --- a quarter of the
signal's standard deviation. Each stretch is normalised separately. \\
7 · does a longer window fix the decomposition? &
\textbf{Partly.} The annual amplitude's overestimate falls from 13.7\,\% to
5.0\,\%, but the trend's share of low-frequency variance \emph{rises} from
39.9\,\% to 53.8\,\%. \fail\ as a general fix: NeuralProphet's trend absorbs the
annual cycle at either length. \\
7 · prediction, full era vs block &
Lagged MAE $-44$\,\% at 1\,h and $-18$\,\% at 168\,h; intervals fall from
164--185\,mdeg at 100\,\% coverage to \textbf{12--23\,mdeg at 55--84\,\%}. Future
regime improves 26--46\,\% at every horizon. \\
7 · the bridged block &
3.71 cycles for \textbf{14.3\,\%} fabricated data, biasing the annual amplitude
$-3.4$\,\% and every shorter component by 15--17\,\%. Reported, not recommended. \\
7 · \texttt{drop\_missing} &
\fail\ --- NeuralProphet 0.8.0's own mechanism for gapped autoregression returns
a NaN loss from the first epoch and cannot predict at all. The multi-series route
is used instead. \\
\bottomrule
\end{tabularx}
\end{center}

% ===========================================================================
\section{Caveats}

\textbf{The target is constructed, and the sensitivity table shows how much that
matters.} Every absolute error depends on a compensation coefficient never
calibrated for this instrument. The relative statements --- which components
matter, which do not --- are far more robust than the levels.

\textbf{The sign contradiction is now two-instrument and still unresolved.}
Section~\ref{sec:sign} adds evidence but does not settle it.

\textbf{The 113.9-day component is recorded, not explained.} It is genuine by the
window-function and permutation controls, and this report offers no mechanism.

\textbf{Non-stationarity is diagnosed, not modelled.} The study establishes that
the annual amplitude and phase move; it does not fit a time-varying seasonal
model, which is the obvious next step and would change the
Section~\ref{sec:budget} conclusions if it succeeded.

\textbf{One station, one channel.} The legacy network had three blocks. This
study reads only \texttt{st02} for comparability with the extended-sensor study,
so it cannot separate structural from instrument behaviour by cross-station
comparison --- the test most likely to resolve Section~\ref{sec:sign}.

\textbf{The 438-day outage cannot be bridged, and the bridged block is not a
recommendation.} Section~\ref{sec:usable} builds it because the question was
asked and because the cost of an answer should be measured rather than asserted.
The holes it fills are far longer than anything the Section~\ref{sec:imputation}
benchmark licenses, and no result in this report depends on it.

\textbf{Sections 1 to 6 keep the longest-block restriction.} They were computed
under it, and they are left that way so that every ported result stays comparable
with the extended-sensor study, which had no choice in the matter.
Section~\ref{sec:usable} measures what the restriction costs; it does not
retrofit the earlier sections with the answer.

% ===========================================================================
\section{Artefact inventory}

Artefacts are numbered in reading order.

\begin{center}
\small
\begin{tabularx}{\textwidth}{L{4.2cm}Y}
\toprule
\textbf{Artefact} & \textbf{Content} \\
\midrule
\texttt{S4\_01} -- \texttt{S4\_04} & Channel summary, coverage, gap inventory,
contiguous blocks \\
\texttt{S4\_05} -- \texttt{S4\_07} & Sign test; correlation matrices, levels and
differences \\
\texttt{S4\_08}, \texttt{S4\_09} & Spectral peaks, low and high band \\
\texttt{S4\_10} -- \texttt{S4\_13} & Harmonic fit, sliding fit, stability,
phase-lock test \\
\texttt{S4\_14}, \texttt{S4\_15} & Imputation benchmark and policy \\
\texttt{S4\_16} -- \texttt{S4\_19} & Operators full band, diurnal band, signed;
phase control \\
\texttt{S4\_20} & Variance inflation factors \\
\texttt{S4\_21}, \texttt{S4\_22} & Single-predictor rankings, whole era and
unbroken block \\
\texttt{S4\_23} & Exhaustive subset search \\
\texttt{S4\_24} -- \texttt{S4\_26} & Leave-one-out gain, permutation importance,
out-of-period test \\
\texttt{S4\_27} & The diagnostic protocol \\
\texttt{S4\_28} -- \texttt{S4\_31} & Error budget, marginal gains, trend cost,
regime gap \\
\texttt{S4\_32} -- \texttt{S4\_34} & NeuralProphet horizons, limits, joint head \\
\texttt{S4\_35}, \texttt{S4\_36} & Sensitivity sweep, verdicts \\
\texttt{S4\_37}, \texttt{S4\_38} & Gap-tolerance sweep; window accounting per
fill policy \\
\texttt{S4\_39}, \texttt{S4\_40} & Baseline step across the long outage;
decomposition, control against the whole era \\
\texttt{S4\_41} -- \texttt{S4\_43} & Segmented horizons, gain over the block,
bridged-block bias \\
\addlinespace
\texttt{S4\_F01}, \texttt{S4\_F02} & Overview and target \\
\texttt{S4\_F03}, \texttt{S4\_F04} & Spectra, low and high band \\
\texttt{S4\_F05} -- \texttt{S4\_F07} & Harmonic fit, sliding harmonics, diurnal
cycle \\
\texttt{S4\_F08}, \texttt{S4\_F09} & Gaps and imputation; correlations \\
\texttt{S4\_F10} -- \texttt{S4\_F12} & Operator grids: full band, diurnal,
signed \\
\texttt{S4\_F13} -- \texttt{S4\_F15} & Ranking, subsets, contribution
diagnostics \\
\texttt{S4\_F16} & Error budget \\
\texttt{S4\_F17} -- \texttt{S4\_F19} & NeuralProphet skill, calibration, fan
chart \\
\texttt{S4\_F20} & Sensitivity \\
\texttt{S4\_F21} -- \texttt{S4\_F23} & Window accounting; whole-era
decomposition; segmented skill \\
\bottomrule
\end{tabularx}
\end{center}

Reproduction, from the study directory with \texttt{neuralprophet\_env} active:

\begin{quote}
\ttfamily\small
jupytext -{}-to ipynb inclination\_prediction\_legacy\_study.py \\
jupyter nbconvert -{}-to notebook -{}-execute -{}-inplace \\
\phantom{jupyter nbconvert }inclination\_prediction\_legacy\_study.ipynb
\end{quote}

\end{document}
